\documentclass[letterpaper,9pt]{article}
\usepackage[utf8]{inputenc}
\usepackage[T1]{fontenc}
\usepackage{latexsym}
\usepackage[empty]{fullpage}
\usepackage{titlesec}
\usepackage{marvosym}
\usepackage[usenames,dvipsnames]{color}
\usepackage{xcolor}
\usepackage{enumitem}
\usepackage[colorlinks=true]{hyperref}
\usepackage{fancyhdr}
\usepackage{tabularx}
\usepackage{iftex}
\ifPDFTeX
  \input{glyphtounicode}
  \pdfgentounicode=1
  \usepackage{microtype}
  \DisableLigatures[f]{encoding = *, family = *}
\else
  \usepackage{fontspec}
  \setmainfont{lmroman10-regular.otf}[
    BoldFont=lmroman10-bold.otf,
    ItalicFont=lmroman10-italic.otf,
    BoldItalicFont=lmroman10-bolditalic.otf,
    SmallCapsFont=lmromancaps10-regular.otf,
    Ligatures=NoCommon]
\fi

% Blue accent matching the Word version (#2563EB)
\definecolor{accent}{HTML}{2563EB}
\definecolor{navy}{HTML}{1F2937}
\definecolor{greytext}{HTML}{555555}

\hypersetup{
  linkcolor=accent,
  urlcolor=accent,
  citecolor=accent,
}

\pagestyle{fancy}
\fancyhf{}
\renewcommand{\headrulewidth}{0pt}
\renewcommand{\footrulewidth}{0pt}

% Margins
\addtolength{\oddsidemargin}{-0.55in}
\addtolength{\evensidemargin}{-0.55in}
\addtolength{\textwidth}{1.1in}
\addtolength{\topmargin}{-0.6in}
\addtolength{\textheight}{1.3in}

\urlstyle{same}
\raggedright
\raggedbottom
\setlength{\tabcolsep}{0in}

% Section titles — blue, matching Word headers
\titleformat{\section}{
  \vspace{-1pt}\color{accent}\scshape\raggedright\large
}{}{0em}{}[\color{accent}\titlerule \vspace{-1pt}]
\titlespacing*{\section}{0pt}{6pt}{4pt}

% Custom commands
\newcommand{\resumeItem}[1]{\item\small{#1}}
\newcommand{\resumeSubheading}[4]{
  \vspace{-1pt}\item
  \begin{tabular*}{0.97\textwidth}[t]{l@{\extracolsep{\fill}}r}
    \textbf{#1} & \textit{\color{greytext}#2} \\
    \textit{\small\color{greytext} #3} & \textit{\small\color{greytext} #4}
  \end{tabular*}\vspace{-2pt}
}
\newcommand{\resumeProjectHeading}[2]{
  \item
  \begin{tabular*}{0.97\textwidth}{l@{\extracolsep{\fill}}r}
    \small #1 & #2
  \end{tabular*}\vspace{-1pt}
}

\newcommand{\resumeSubHeadingListStart}{\begin{itemize}[leftmargin=0.18in, label={}, itemsep=1.2pt, topsep=2.5pt, parsep=0pt]}
\newcommand{\resumeSubHeadingListEnd}{\end{itemize}\vspace{-2pt}}

\newcommand{\resumeItemListStart}{\begin{itemize}[leftmargin=0.24in, label={}, itemsep=0.9pt, topsep=1.2pt, parsep=0pt]}
\newcommand{\resumeItemListEnd}{\end{itemize}\vspace{-2pt}}

\newcommand{\resumeProjectListStart}{\begin{itemize}[leftmargin=0.24in, label={}, itemsep=0.9pt, topsep=1.2pt, parsep=0pt]}
\newcommand{\resumeProjectListEnd}{\end{itemize}\vspace{-2pt}}

\newcommand{\resumeLangHeading}[2]{
  \vspace{-1pt}\item
  \begin{tabular*}{0.97\textwidth}[t]{l@{\extracolsep{\fill}}r}
    \textbf{#1} & #2
  \end{tabular*}\vspace{-1pt}
}

\begin{document}

% ---------- HEADER ----------
\begin{center}
    \textbf{\Huge \color{navy} Mohamed Amine Nasfi} \\[3pt]
    \small \textbf{\color{accent} Technicien Support Systèmes \& Réseaux | Postes de travail \textperiodcentered{} Parc et inventaire \textperiodcentered{} Support N1 \textperiodcentered{} Réseaux locaux} \\[3pt]
    \footnotesize \color{greytext}
    +216 27 711 810 $\quad|\quad$
    \href{mailto:medaminenasfy@gmail.com}{\underline{medaminenasfy@gmail.com}} $\quad|\quad$
    Gabès, Tunisie — mobilité nationale \\[2pt]
    \href{https://www.linkedin.com/in/mohamed-amine-nasfi/}{\underline{LinkedIn}} $\quad|\quad$ \href{https://github.com/medaminenasfi}{\underline{GitHub}} $\quad|\quad$ \href{https://mohamedamine-portfolio.onrender.com}{\underline{Portfolio}}
\end{center}
\vspace{-5pt}

% ---------- SUMMARY ----------
\section{Résumé}
\begin{itemize}[leftmargin=0.18in, label={}, itemsep=1.2pt, topsep=1.5pt, parsep=0pt]
\item \footnotesize{Technicien informatique en support de premier et deuxième niveau auprès d'utilisateurs non-techniques : installation, configuration et maintenance des postes de travail Windows 10/11 et des périphériques, diagnostic et résolution des incidents matériels, logiciels et réseau sous délai contraint. Assure le suivi du parc, la mise à jour de l'inventaire, l'enregistrement des demandes et le compte rendu des interventions. Intervient aussi sur l'infrastructure réseau locale : adressage IPv4, DHCP/DNS, configuration routeur et switchs managés (VLAN), pare-feu, sécurisation Wi-Fi et accès distant VPN/SSL, câblage Ethernet. Administre des serveurs Linux en production (VPS, Docker, Nginx, suivi des logs). Licence en Systèmes d'Information — Mention Très Bien, français courant (TCF B2) et anglais professionnel.}
\end{itemize}\vspace{-3pt}

% ---------- EXPERIENCE ----------
\section{Expérience}
\resumeSubHeadingListStart

  \resumeSubheading{Technicien Support Informatique, Réseaux \& Parc --- Freelance}{Mai 2026 -- Septembre 2026}{Cabinet d'Orthodontie Dr Meriem Nasfi}{Gabès, Tunisie}
    \resumeItemListStart
      \resumeItem{Support technique de premier et deuxième niveau auprès d'utilisateurs non-techniques : qualification et résolution des incidents matériels et logiciels sous 2 heures, assistance à distance, suivi des demandes jusqu'à résolution et compte rendu écrit des interventions.}
      \resumeItem{Installation, configuration et maintenance des postes de travail Windows 10/11 et des périphériques : formatage et masterisation, mise en service d'imprimantes, remplacement de composants, maintenance préventive du parc et inventaire tenu à jour.}
      \resumeItem{Déploiement de l'infrastructure réseau locale : plan d'adressage IPv4, configuration routeur et switchs, pare-feu, sécurisation Wi-Fi et accès distant VPN/SSL, tirage et raccordement du câblage Ethernet, vérification de la connectivité et du bon fonctionnement des équipements.}
    \resumeItemListEnd

  \resumeSubheading{Administrateur Systèmes Linux \& Développeur d'Outils --- Freelance}{Nov 2025 -- Présent}{Indépendant (Freelance / Remote)}{Remote}
    \resumeItemListStart
      \resumeItem{Administration et sécurisation de serveurs Linux en production : hébergement VPS, conteneurisation Docker multi-services, reverse-proxy Nginx, renouvellement automatisé des certificats SSL, surveillance de la disponibilité et analyse des logs.}
      \resumeItem{Diagnostic et résolution d'incidents d'exploitation, reprise de service après incident, déploiements automatisés via pipelines CI/CD GitHub Actions et outils internes sur bases PostgreSQL.}
    \resumeItemListEnd

  \resumeSubheading{Développeur Full-Stack --- Projet de Fin d'Études (PFE)}{Fév 2025 -- Juin 2025}{Laghazala du Désert}{Gabès, Tunisie}
    \resumeItemListStart
      \resumeItem{Développement d'une plateforme full-stack (React.js, Node.js, Express.js, MongoDB) avec authentification JWT et chif\hspace{0pt}frement SSL ; recueil des besoins, documentation technique et coordination à distance avec l'encadrant.}
    \resumeItemListEnd

  \resumeSubheading{Stagiaire --- Outils internes \& Documentation}{Juil 2024 -- Août 2024}{Groupe Chimique Tunisien (GCT)}{Gabès, Tunisie}
    \resumeItemListStart
      \resumeItem{Développement d'un outil d'automatisation de données (PHP, SQL) réduisant le temps de saisie de \textbf{70\%} pour une équipe opérationnelle de 5 personnes.}
      \resumeItem{Recueil des besoins auprès des utilisateurs métiers, rédaction des procédures techniques et du guide d'utilisation.}
    \resumeItemListEnd

\resumeSubHeadingListEnd\vspace{-2pt}

% ---------- TECHNICAL SKILLS ----------
\section{Compétences Techniques}
\begin{itemize}[leftmargin=0.18in, label={}, topsep=2pt, itemsep=1pt, parsep=0pt]
\footnotesize{
  \item \textbf{Support utilisateurs N1/N2:} Réception et qualification des demandes, résolution des incidents matériel et logiciel, assistance sur site et à distance (AnyDesk, TeamViewer), suivi des demandes jusqu'à résolution, comptes rendus écrits, formation des utilisateurs non-techniques.
  \item \textbf{Postes de travail \& Matériel:} Windows 10/11 : installation, formatage, masterisation, pilotes, registre et dépannage ; montage et mise en service de postes ; remplacement de composants (RAM, SSD, alimentation) ; imprimantes et périphériques ; maintenance préventive et corrective.
  \item \textbf{Parc informatique \& Inventaire:} Inventaire du parc tenu à jour (tableur), suivi des entrées et sorties de matériel, vérification du bon fonctionnement des équipements, standardisation des configurations et checklists d'installation, documentation des procédures d'intervention.
  \item \textbf{Réseaux:} TCP/IP, adressage IPv4 et sous-réseaux, DHCP/DNS, routeurs et switchs managés (VLAN), pare-feu et segmentation, redirection de ports, sécurisation Wi-Fi, accès distant VPN/SSL, diagnostic des pannes réseau local, câblage Ethernet RJ45 (sertissage, test).
  \item \textbf{Systèmes \& Serveurs:} Linux Ubuntu/Debian en production : ligne de commande et bash, hébergement VPS, Docker, Nginx, certificats SSL, suivi des logs et de la disponibilité des services, sauvegardes ; exploitation des postes et services Windows en environnement professionnel.
  \item \textbf{Bureautique \& Outils:} Pack Of\hspace{0pt}fice / Microsoft 365 en usage utilisateur (traitement de texte, tableur, messagerie, réunion à distance) ; Git/GitHub, Postman ; automatisation de tâches en Python et Bash ; bases de données PostgreSQL, MySQL, SQLite.
}
\end{itemize}\vspace{-3pt}

% ---------- PROJECTS ----------
\section{Projets}
\resumeSubHeadingListStart

  \resumeProjectHeading{\textbf{Application de gestion et facturation --- Bureau Windows} $\quad$ --- $\quad$ \emph{\color{greytext}Python \textperiodcentered{} PyQt5 \textperiodcentered{} SQLite \textperiodcentered{} PyInstaller}}{}
    \resumeProjectListStart
      \resumeItem{Application de bureau déployée sur postes Windows pour une activité de location : contrats, facturation, tableau de bord et sauvegardes automatisées horodatées de la base SQLite.}
    \resumeProjectListEnd

\resumeSubHeadingListEnd\vspace{-3pt}

% ---------- EDUCATION ----------
\section{Formation}
\resumeSubHeadingListStart
  \resumeSubheading{Licence en Systèmes d'Information --- Mention Très Bien}{2022 -- 2025}{ESSAT Gabès}{Gabès, Tunisie}
\resumeSubHeadingListEnd\vspace{-1pt}

% ---------- LANGUAGES ----------
\section{Langues}
\begin{itemize}[leftmargin=0.18in, label={}, topsep=2pt, itemsep=1pt, parsep=0pt]
\footnotesize{
  \item Français (Courant --- B2 (TCF TP)) \textperiodcentered{} Anglais (Professionnel) \textperiodcentered{} Arabe (Langue maternelle)
}
\end{itemize}

\end{document}
